% !TeX root = ../papers/nc-p3.tex
\section{滑らかな二次曲面のグラフが続く窓の全分類}
\label{sec:smooth-quadric-graphs}
四辺の各射影の非空なファイバーがすべて零次元であり、少なくとも一辺の二文字点対応が
スキームとして滑らかな二次曲面間の同型のグラフであるとする。
他の辺の射影像の一致や全位置を通る点列は仮定しない。
\begin{theorem}\label{thm:smooth-quadric-graphs}
この枝の標準14641型の添字保持同型類は $\mathbb A^1\setminus\{2\}$ である。
四文字 $x,y,z,w$ を、二平面 $\langle x,y\rangle,\langle z,w\rangle$ に分ける。
異なる平面の二文字は交換し、ほかの二関係は位置 $i$ で
\[
 xy+yx+zw+wz=0,\qquad [z,w]=(-1)^i m[x,y].
\]
パラメータは $m\ne0,\ m^2\ne1$ で、上のアフィン座標は $m^2+m^{-2}$ である。
全点の最小周期は1である。これらの窓は正則正規一次直線列を持たず、
どの量子多項式点も軌道閉包に持たない。
\end{theorem}
\begin{proof}[二次核の正規性]
まず全辺のグラフがつながることを示す。各射影のファイバーは線型なので、
零次元条件は左右評価の全点での階数3以上に等しい。辺上の $(p,r)$ では
命題\ref{prop:three-finite-edges}の一般不等式 $b_i(p)+f_{i+1}(r)\le6$ により両方3である。
系\ref{cor:finite-edge-point-sequences}の局所証明で必要なのはこの階数だけである。
すなわち基点で単元な三次小行列を分裂させると、次辺評価の四次小行列がスキーム上で消え、
その核直線への射が定まる。逆向きの核直線も一意であるから共有文字への両像はスキームとして一致する。
これを反復すれば四辺とも滑らかな二次曲面のグラフである。
二次曲面の同型は二直線族を保つか交換するため、超平面束 $\mathcal O(1,1)$ を保つ。
従って点対応に沿って全整数位置の基底を運び、一つの $Q\subset\mathbb P(V^*)$ の
対角グラフにできる。$q\in\operatorname{Sym}^2 V$ を非退化なその方程式の対称持上げとする。
グラフを消す二次式の空間は $K=\bigwedge^2 V+kq$ である。各六関係はこの中の超平面で、
一次元像の生成元 $s_i$ を選ぶと
\[
 [p,r]_i=\Omega_i(p,r)s_i,\qquad q_i=c_i s_i
\]
と書ける。$\Omega_i=0$ なら辺は平面三空間全体の対角である。
階数2なら $\Omega_i$ を $V^*$ の可分解二ベクトルと見て、その二平面の射影直線と
$Q$ の交点 $[p]$ を取る。$(p,p)$ の接条件の交代部分は
$p\wedge(\dot p_2-\dot p_1)\in k\Omega_i$ で、差には射影接方向が一つ残る。
二つの変位を合わせた次元4から残る $q$ の一式を引いても接次元は3以上であり、
滑らかな対角グラフに反する。よって全ての $\Omega_i$ は非退化である。
逆に非退化なら $p\wedge r=t\Omega_i$ に $p$ を外積すると $p\wedge\Omega_i\ne0$ から
局所的に $t=0$ となり、残る条件は対角と $q(p,p)=0$ そのものである。

全 $s_i$ による商は可換二次曲面環 $D=k[V]/(q)$ の反復である。
$\ell_\Omega:\bigwedge^3 V\to V$ を
$\ell_\Omega(u,v,h)=\Omega(u,v)h+\Omega(v,h)u+\Omega(h,u)v$ とする。
三文字の交代和から
$s_i\ell_{\Omega_i}(u,v,h)=\ell_{\Omega_{i+1}}(u,v,h)s_{i+1}$ である。
両 $\ell$ は可逆なので $s_i V=V s_{i+1}$ である。これを全経路に繰り返せば
商への核はそれぞれ片側の $s$ の倍で尽くされる。標準Hilbert関数と
$\dim D_n=(n+1)^2$ の差は $\dim A_{i+2,i+n}$ に等しいから、左右の $s$ 乗法は全て単射である。
従って $s_i v=g_i(v)s_{i+1}$ は二移動の可逆な乗法移送を与え、さらに $s$ の
イデアルを保つ（$s_i s_{i+2}$ の式で両端の正則性を使う）。
商 $D$ の点対応は $Q$ の対角だから $g_i$ は全位置で射影的に同じで、
その射影変換は $Q$ を保つ。$s_i$ の倍率を順に直して $g_i=g$ と取る。
\end{proof}

\begin{proof}[二直交平面の復元]
$q$ の対称行列を $J:V^*\to V$ とし、$gJg^t=\mu J$ と書く。
体積を用いて $\ell_{\Omega_i}$ を可逆交代行列 $M_i:V^*\to V$ と見る。
共通の符号規約の下で $M_i$ は $\operatorname{Pf}(\Omega_i)\Omega_i^{-1}$ の
一定の非零倍である。三文字の交代恒等式から $M_{i+1}=g M_i$、
ゆえに $g^t\Omega_i=\Omega_i g$ となる。$s=\operatorname{Pf}(M_{i+1})/\operatorname{Pf}(M_i)$ とすると
$\Omega_{i+1}=s\Omega_i g^{-1}$。各一般化固有空間は交代形式で互いに直交し非退化なので
各固有値の重複度は偶数である。この比 $s$ は各固有値をその半重複度だけ掛けた値である。
実際 $\operatorname{Pf}((T I-g)M_i)/\operatorname{Pf}(M_i)$ は
$\det(T I-g)$ のモニック平方根で、その定数項がこの比である。
従って $s$ は位置によらない。

$qv-vq$ を二つの二文字交換子に展開すると
\[
 g(c_i I-J\Omega_i)=c_{i+1} I+J\Omega_{i+1}.
\]
ここで $\Omega_i$ の行列要素は $\Omega_i(e_j,v)$ の規約で取る。
相似式と自己随伴性から $gJ\Omega_i=(\mu/s)J\Omega_{i+1}$ なので
\[
 c_i g-c_{i+1} I=(1+\mu/s)J\Omega_{i+1}.
\]
$g$ がスカラーなら左辺はスカラー、右辺は非零な反対称形式に由来する非スカラーであり不可能である
（このとき $\mu=s$）。$\mu\ne -s$ なら $J\Omega_{i+1}$ は $g$ と可換する。
$J\Omega_i=s^{-1}J\Omega_{i+1}g$ も可換し、$gJ\Omega_i=\mu J\Omega_i g^{-1}$ から
$g^2=\mu I$ である。非スカラーなので二固有値は $b,-b$ で重複度は2ずつ、
$s=-b^2=-\mu$ となって矛盾する。よって $\mu=-s$ で全 $c_i=0$ である。

固有値が $b$ 一つなら、対称形式の相似性から $\mu=b^2=s$ が必要であり不可能である。
二固有値 $b,c$ の場合 $s=bc$ であり、$\mu=-bc$ では両一般化固有空間は
対称形式でも互いに直交する。それぞれで非退化なので $b^2=c^2=\mu$、従って $c=-b$。
二次元の非退化交代平面上の自己随伴写像はスカラーであり、Jordan部分はない。
$g$ の全体倍率を直して二固有値を $1,-1$ と取る。
両平面で $q$ を双曲型に取れば $q=xy+yx+zw+wz$、二平面間の交換子は零である。
$\Omega_{i+1}=s\Omega_i g^{-1}$ の二平面上の倍率の比が $-1$ なので、
超平面の式は定理の表示で尽くされる。
\end{proof}

\begin{proof}[双対Clifford環と除外される二値]
$D^!$ を通常次数で書いた環を $C$ とする。双対文字を $X,Y,Z,W$ とすると
四文字は平方零、二平面間は反交換、$XY+YX=ZW+WZ=h$ は中心である。
$C$ は $k[h]$ 上で十六の平方自由順序語を基底とする。実際各二文字対を
$YX=h-XY$ 等で並べ、平方と混合する重なり $YXX,YYX$ は零、二平面は符号付きで交換する。
これで $h$ を独立変数とした表示の基底が得られ、$h$ を消すと元の二次表示になる。
$U=XY-YX$, $T=ZW-WZ$ と置くと $U^2=T^2=h^2$, $UT=TU$ である。
$A^!$ は各位置で $C$ の一関係 $a_i=U+(-1)^i mT$ を加えたものとなる。
実際この関係の $K$ 上の汎関数は $q$ で零、二交換子の上で比 $(-1)^i m$ である。
$U$ は $X,Y$ と反交換し $Z,W$ と交換、$T$ はその逆なので $a_i$ は次位置の
$a_{i+1}$ へ符号付きで移り、$a_i a_{i+1}=(1-m^2)h^2$ である。

$m^2=1$ なら $a_{i+1}h^n\ne0$ が全て $a_i$ で消える。
$C_j$ は $j\ge3$ で常に次元8なので、正規性により商の高次偶数片は無限に残る。
標準型ではない。$m^2\ne1$ なら $a_i$ は左右正則で、商の各始点級数は $(1+z)^4$ である。
可換多項式環の正則二次式商 $D$ はKoszulであり、その双対 $C$ もKoszulである。
$C$ のこの正則正規二次式商もKoszulである。実際、変換環Tor系列の二段は
商のTorの下段と内部次数2を移した上段である。最初の非線型下段成分の出る微分は
二列前の線型な上段に入るので非線型次数を消せず、元のKoszul性に反する。
位置別の符号移送は頂点や矢の正則な捩りだけを加え、次数比較は同じである。
同じ議論を通常多項式環から $D$ にも適用できる。

$F=C/(h^2)$ は最高次数6のFrobenius環である。socleは $h$ に消されるので $hF$ に入り、
$C/(h)$ が四文字の外積環であることから左も右も $hXYZW$ の一本だけである。
位置 $i$ での商 $C/(a_i)$ の次数片を、$a_{i+1}$ の左乗法で $F$ の次数を2上げた部分へ入れる。
これは $a_i a_{i+1}=(1-m^2)h^2$ と単射性から核が正確に $a_i C$ である。
$F$ の積と最高traceを使えば、商の各非零な次数 $j$ の元に、右から掛けて非零となる
次数 $4-j$ の元がある。$a$ の正規性によりこの積は相手側でも商を通る。
相補次元が等しいので全積対合が完全である。よって $A^!$ はKoszulかつ標準Frobenius、
元の窓は全て標準Koszul左右AS正則である。
\end{proof}

\begin{proof}[同型と二つの補集合への所属]
点グラフから正規化後の全位置の同型は射影的に同じで $Q$ を保つ。
復元した $g$ の二固有平面を保つか交換する。保つものは二平面で同じ相似倍率を持ち、
交代面積に掛かる二つの行列式の比は $\pm1$ である。両符号が実現し、平面交換は比を逆にする。
ゆえに同型条件は $m'\in\{m,-m,m^{-1},-m^{-1}\}$ である。
この作用の不変環は $k[m^2+m^{-2}]$、除外する値は2である。一位置移動は $m\mapsto-m$ で同型である。
正規一次列があれば $D$ 上で連続する二つの超平面切断の零因子が等しい。
$Q$ は整で張るので全位置で同じ文字直線 $kv$ になる。$v V=V v$ の対応は、
可換整域 $D$ へ移すと各文字を固定するから、元の $A$ でも $[v,V]=0$ を要する。
$\Omega_i$ の非退化性に反する。

最後にこの窓の射影軌道はpolystableである。上で得た最高次四テンソルは、各語について
$F$ の $hXYZW$ 係数を $a_{i+1}$ とその四文字の積から取って表示できる。
$(X,Y,Z,W)$ を $(tX,t^{-1}Y,uZ,u^{-1}W)$ に動かすトーラスで不変なので、
異なる一文字切片の台は交わらない。語の台は $(XY)(ZW)$ の同時入替えで保たれる。
また一位置巡回と $Z,W$ 入替えの合成でも保たれる。実際、$F$ の最高traceに関する
一文字の巡回は符号 $-1$ を掛ける。平方自由基底で、相手の次数5が $h$ 倍の三文字であることから
外積の符号となる。さらに $a_{i+1}$ を一文字越すと符号を除き $a_i$ に変わり、
$Z,W$ 入替えがこれを元に戻す。この有限置換群で各語を平均すると、
各位置で第一平面の文字の重みは二文字とも同じで、他平面も同様、各平面の総重みは位置によらない。
台には $XYXY$ と $ZWZW$ がともにある（係数はそれぞれ $m$ および1の非零倍）。
従って全語に正の重みを与えた平均にこの二軌道の重みを調節し、全周辺重みを均等にできる。
相対内部の重心による対角ノルム最小化と切片の直交性
（定理\ref{thm:quantum-face-resonance}の証明）から全四基底群でpolystableとなる。
量子点の射影支持は座標線型空間の和であり、この滑らかな二次曲面とは異なる。
異なるpolystable軌道は互いの閉包に入らない。
\end{proof}
